\documentclass[11pt,a4paper]{article}

\usepackage[T1]{fontenc}
\usepackage[utf8]{inputenc}
\usepackage{lmodern}
\usepackage[a4paper,margin=26mm]{geometry}
\usepackage{amsmath,amssymb,amsthm,mathtools}
\usepackage{booktabs,array}
\usepackage{microtype}
\usepackage{xcolor}
\usepackage{enumitem}
\usepackage{url}
\usepackage[colorlinks=true,linkcolor=blue!55!black,citecolor=blue!55!black,
            urlcolor=blue!55!black]{hyperref}

\newtheorem{theorem}{Theorem}
\newtheorem{lemma}[theorem]{Lemma}
\newtheorem{proposition}[theorem]{Proposition}
\newtheorem{remark}[theorem]{Remark}
\newcommand{\Z}{\mathbb{Z}}
\newcommand{\R}{\mathcal{R}}
\newcommand{\E}{\mathbb{E}}
\newcommand{\Cov}{\mathop{\rm Cov}}
\newcommand{\round}[1]{\left\lfloor #1\right\rceil}
\newcommand{\norm}[1]{\left\lVert #1\right\rVert}
\newcommand{\pk}{\mathsf{pk}}
\newcommand{\sk}{\mathsf{sk}}
\newcommand{\Sign}{\mathsf{Sign}}
\newcommand{\Verify}{\mathsf{Verify}}
\newcommand{\KeyGen}{\mathsf{KeyGen}}
\newcommand{\Compress}{\mathsf{CompressY}}
\newcommand{\Stretch}{\mathsf{StretchS}}

\title{\bfseries Covariance Leakage in the NGCC Shuttle Submission\\
\large Equivalent-Key Recovery and Fresh-Message Forgery\\
Against All Three Submitted Parameter Sets}
\author{Martin Feussner\\
\small Selmer Center, University of Bergen\\
\small \texttt{martin.feussner@uib.no}}
\date{25 September 2026}

\begin{document}
\maketitle

\begin{center}
\fbox{\parbox{0.92\linewidth}{\small\textbf{Provenance and review status.}
The attack, derivation, implementation, experiments, and initial manuscript
were produced independently by OpenAI Codex using Daybreak Blue at max
reasoning effort during an AI-run audit of NGCC signature candidates.  At the
date above, no human cryptanalyst has independently verified the argument,
implementation, measurements, or novelty assessment.  This report is a
preliminary disclosure for human review and reproduction.}}
\end{center}
\vspace{0.7em}

\begin{abstract}
Shuttle is a lattice signature submitted to the Next-generation Commercial
Cryptographic Algorithms Program.  Its rejection-free iterative sampler is
intended to select between $y-v$ and $y+v$ so that the response remains close
to a secret-independent Gaussian.  The submitted specification defines
$p_v(y)$ as the probability of $y-v$, but its transition pseudocode applies
the same event to $y+v$.  The submitted reference implementation follows this
reversed update.  The authors' independently published ePrint version uses
the opposite, intended branch direction.  This establishes a concrete mismatch
between the NGCC submission and the separately published construction; the
attack described here targets that NGCC-specific mismatch.

The reversal changes a hiding transition into an amplifying one.  In the
idealized exact-transition analysis, reversing the branch changes covariance
cancellation into an additional rank-one term $2vv^T$.  Shuttle's first secret component is the
public constant one, so the transmitted first response polynomial supplies an
anchor for every nonzero challenge position.  Cross-correlating that anchor
with the other transmitted response polynomials estimates every coefficient
of the signing secret directly; after Shuttle's matching stretch and
compression factors cancel, the signal is approximately $2s_i$.

We give a streaming chosen-message attack using only ordinary valid
signatures.  Rounding the empirical covariances recovers $s_1$.  The omitted
error vector is then computed exactly from the public relation
$b=a_{\rm gen}+A_{\rm gen}s_1+s'_2\pmod q$.  Public coefficient bounds, the
key-generation norm window, and the public key provide an exact validation
test.  The resulting equivalent secret key signs arbitrary messages.

End-to-end experiments used the unmodified submitted reference signer and
verifier.  Shuttle-128 was broken for four independent keys using
150,000--225,000 signatures and 311--467 single-threaded CPU-seconds.
Shuttle-256 was broken with 225,000 signatures in 740 CPU-seconds.
Shuttle-512 was broken with 325,000 signatures in 2,208.998
CPU-seconds.  Every queried signature was verified, and each recovered key
produced an accepted signature on a fresh message.  Peak resident memory was
below 2.5 MB.  All runs used a 3.2 GHz Intel Core i5-6500 with 16 GiB
of RAM, placing the complete attack well within commodity desktop resources.
\end{abstract}

\section{Introduction}

Fiat--Shamir lattice signatures add a secret-dependent shift to a wide random
mask.  Their rejection or transition rule has one central privacy obligation:
the public response distribution must be independent of the secret.  A sign
error in that rule is especially dangerous because honest signatures then
become repeated noisy measurements of the signing key.

Shuttle is an NGCC lattice-signature submission based on MLWE and MSIS
\cite{Shuttle}.  It adapts iterative rejection sampling \cite{Gaertner} into a
two-output transition with no inner abort.  A sparse challenge selects cyclic
shifts of a stretched secret vector, and the transition chooses the sign of
each shift.  The submission claims that the accumulated response is close to
a public discrete Gaussian, with the remaining approximation controlled by
R\'enyi divergence.

The submitted transition does the opposite of its own probability
definition.  In the submission's notation, $p_v(y)$ is the probability of
returning $y-v$.  Its interval test recognizes an event of probability
$p_v(y)$, sets a flag to $+1$ on that event, and returns $y+\mathit{flag}\,v$.
The portable C implementation makes the identical update.  Thus it returns
$y+v$ where the analysis requires $y-v$, and vice versa.  The construction published separately by the Shuttle authors explicitly
returns $y-v$ on the interval event \cite{MingEtAl}.  This provides a direct
cross-check of the branch convention used by the NGCC submission.

The resulting leak is visible in second moments.  The correct transition must
cancel the covariance introduced by adding a signed $v$.  Reversing its
output sign doubles that covariance instead.  Shuttle makes this exploitable
without learning the sampler's internal signs: the full secret is
$(1,s_1,s'_2)$, and the response block corresponding to the known component
$1$ is transmitted.  It exposes a noisy copy of each internal sign at the
corresponding challenge position.  The other transmitted blocks then expose
the aligned coefficients of $s_1$.

Our contributions are:

\begin{itemize}[leftmargin=1.5em]
  \item We identify the reversed branch in both the NGCC specification and
        the submitted portable implementation, and contrast it with the
        correctly oriented transition in the authors' public ePrint.
  \item We derive the rank-one covariance term caused by the reversal and
        turn it into a coefficient-wise estimator over
        $\Z[X]/(X^n+1)$.
  \item We give a streaming recovery algorithm that stores only $3n$
        signed accumulators and uses no signing randomness, side channel,
        malformed input, or hash weakness.
  \item We complete the recovered $s_1$ to an equivalent signing key using
        the public key alone and validate it with all public key-generation
        bounds.
  \item We demonstrate accepted fresh-message forgeries against all three
        submitted parameter sets on an ordinary four-core 2015 desktop CPU,
        using far fewer than one million signatures.
\end{itemize}

The attack is on the version submitted to NGCC.  It does not apply to an
implementation of the correctly oriented transition described in
\cite{MingEtAl}.

\section{Relevant parts of Shuttle}
\label{sec:scheme}

Let
\[
       \R_q=\Z_q[X]/(X^n+1).
\]
The public key contains a seed defining
$a_{\rm gen}\in\R_q^m$ and
$A_{\rm gen}\in\R_q^{m\times\ell}$, together with
$b\in\R_q^m$.  Key generation samples
$s_1\in\R^\ell$ and $s_2\in\R^m$, rounds the initial public value, and absorbs
the rounding discrepancy into an adjusted error $s'_2$.  Consequently the
encoded public key satisfies the exact relation
\begin{equation}
       b=a_{\rm gen}+A_{\rm gen}s_1+s'_2\pmod q.
\label{eq:pkrelation}
\end{equation}
The full signing secret used by the iterative sampler is
\[
                         s=(1,s_1,s'_2).
\]

\subsection{Stretching, compression, and the response}

Shuttle uses block-dependent integer factors
$\alpha_1,\alpha_s,\alpha_e$.  Define
\begin{align*}
 \Stretch(x_0,x_s,x_e)
   &=(\alpha_1x_0,\alpha_sx_s,\alpha_ex_e),\\
 \Compress(y_0,y_s,y_e)
   &=\bigl(\round{y_0/\alpha_1},
            \round{y_s/\alpha_s},
            \round{y_e/\alpha_e}\bigr).
\end{align*}
The shifted-floor convention gives the exact translation identity
\begin{equation}
       \Compress\bigl(y+\Stretch(x)\bigr)=\Compress(y)+x
       \qquad(x\text{ integral}).
\label{eq:translation}
\end{equation}

The signer samples a wide discrete-Gaussian vector $y$, computes a Fiat--Shamir
challenge
\[
                  c=\sum_{j\in J}X^j,
                  \qquad |J|=\tau,
\]
and processes the indices $j\in J$ in increasing order.  For each index it
uses the shift
\[
                       v_j=\Stretch(s)X^j
\]
and chooses a hidden $\epsilon_j\in\{-1,+1\}$.  Before compression, the final
state has the form
\[
                \widetilde z=y+\sum_{j\in J}\epsilon_jv_j.
\]
Equation~\eqref{eq:translation} makes the public response exactly
\begin{equation}
       z=\Compress(y)+s\sum_{j\in J}\epsilon_jX^j.
\label{eq:response}
\end{equation}
Only the first $1+\ell$ polynomials, denoted $z_1$, are serialized directly;
the $m$ error-block polynomials are represented through a hint.  The attack
needs only $z_1$ and the public challenge seed.

For every $j\in J$, Equation~\eqref{eq:response} gives
\begin{align}
 z_{1,0}[j] &= \round{y_0[j]/\alpha_1}+\epsilon_j,
 \label{eq:anchor}\\
 z_{1,p}[j+d]_{\rm neg}
   &=\round{y_p[j+d]_{\rm neg}/\alpha_s}
     +\epsilon_js_{1,p}[d]+\text{other shifts},
 \label{eq:secretobs}
\end{align}
where $1\le p\le\ell$, and the signed negacyclic access is
\[
 x[k]_{\rm neg}=\begin{cases}
       x[k],&0\le k<n,\\
      -x[k-n],&n\le k<2n.
 \end{cases}
\]

\subsection{The intended transition}

For a single shift $v$, Shuttle defines
\[
 S_v(y)=\sum_{i=0}^{\infty}(-1)^i
              \frac{\rho_r(y+iv)}{\rho_r(y)}
\]
and
\begin{equation}
 p_v(y)=\begin{cases}
   S_v(y),&\langle y,v\rangle\ge 0,\\
   1-S_v(-y),&\langle y,v\rangle<0.
 \end{cases}
\label{eq:pv}
\end{equation}
Both the NGCC prose and the authors' ePrint define $p_v(y)$ as the probability
of returning $y-v$; $y+v$ has probability $1-p_v(y)$
\cite{Shuttle,MingEtAl}.  The finite alternating sum is implemented by drawing
$U\in(0,1]$ and testing whether $2r^2\log U$ lies in a union of intervals
whose total probability is $p_v(y)$.

After normalizing the sign of $v$ so that $\langle y,v\rangle>0$, let
$B(y,v,U)$ be $+1$ when this interval test succeeds and $-1$ otherwise.
The correct transition is
\begin{equation}
                         y-Bv.
\label{eq:correct}
\end{equation}
Equivalently, the ePrint writes the event bit as $b\in\{0,1\}$ and returns
$y+(-1)^bv$, so membership ($b=1$) returns $y-v$ \cite{MingEtAl}.

\subsection{The submitted sign reversal}
\label{sec:reversal}

Algorithm 22 of the NGCC specification initializes a flag to $-1$, sets it to
$+1$ when the same interval test succeeds, and returns
\begin{equation}
                         y+Bv.
\label{eq:reversed}
\end{equation}
The submitted C implementation also updates every coefficient by
$z\leftarrow z+Bv$.  Thus Equations~\eqref{eq:correct} and
\eqref{eq:reversed} are exact opposites for the same random event.

This is not a harmless renaming of $v$.  The probability in
Equation~\eqref{eq:pv} depends on the oriented inner product, and its purpose
is to counteract the selected shift.  The submitted update reinforces it.

\section{Why the reversal creates covariance}
\label{sec:covariance}

The leak can be seen without evaluating the alternating series.  Consider one
transition with input $Y$ of mean zero and covariance $\Sigma$.  Let
$E\in\{-1,+1\}$ include both the sign normalization of $v$ and the outcome of
the interval test.  The intended output is
\[
                         Z_-=Y-Ev,
\]
whereas the submitted output for the same event is
\[
                         Z_+=Y+Ev.
\]
Write $\mu=\E[EY]$.  Since $E^2=1$,
\begin{align}
 \E[Z_-Z_-^T]
    &=\Sigma-\mu v^T-v\mu^T+vv^T,\label{eq:minusmoment}\\
 \E[Z_+Z_+^T]
    &=\Sigma+\mu v^T+v\mu^T+vv^T.\label{eq:plusmoment}
\end{align}

The intended transition is designed to preserve the Gaussian up to a tiny
pointwise error.  In the idealized exact case,
$\E[Z_-Z_-^T]=\Sigma$.  The resulting matrix identity
$\mu v^T+v\mu^T=vv^T$ itself forces $\mu=v/2$: multiplication by any vector
orthogonal to $v$ first shows that $\mu$ is parallel to $v$, and comparison
along $v$ fixes the scalar.  Substitution into
Equation~\eqref{eq:plusmoment} yields
\begin{equation}
                     \E[Z_+Z_+^T]=\Sigma+2vv^T.
\label{eq:rankone}
\end{equation}
For the approximate discrete transition, the same calculation has an error
controlled by the deviation of the intended output from the target Gaussian.
The submission claims that deviation is negligible at its parameters; it
cannot cancel the leading $2vv^T$ term.

For $v_j=\Stretch(s)X^j$, the first block of $v_j$ has the single nonzero
coefficient $\alpha_1$ at position $j$.  Its cross-covariance with coefficient
$j+d$ of secret block $p$ is therefore
\[
              2\alpha_1\alpha_s s_{1,p}[d].
\]
The public compression divides the two blocks by the same factors.  Thus the
observable cross-covariance in Equations~\eqref{eq:anchor} and
\eqref{eq:secretobs} is
\begin{equation}
 \E\!\left[z_{1,0}[j]\,
       z_{1,p}[j+d]_{\rm neg}\mid j\in J\right]
            \approx 2s_{1,p}[d].
\label{eq:observable}
\end{equation}
Equation~\eqref{eq:observable} motivates the estimator used below.  In the
complete signing process, other challenge shifts, compression and rounding,
later transitions, and the final public norm gate introduce additional terms.
We expect cross-terms associated with unrelated random challenge locations to
average toward zero; rather than assuming that this approximation is exact,
Section~\ref{sec:experiments} measures the resulting coefficient scale
directly.  Across all tested parameter sets the empirical multiplier remains
close to the predicted value $2$, while the corrected-sign control of
Section~\ref{sec:corrected-control} removes the signal.

\section{Key recovery and forgery}
\label{sec:attack}

The complete attack has four stages: collect signatures, estimate $s_1$,
derive $s'_2$ from the public key, and serialize an equivalent signing key.

\subsection{Streaming covariance estimator}

Initialize signed 64-bit accumulators $C_{p,d}=0$ for
$1\le p\le\ell$ and $0\le d<n$, and an observation counter $M=0$.
For each ordinary signature:

\begin{enumerate}[leftmargin=2em]
  \item Verify it with the submitted public verifier and discard it on any
        failure.
  \item Decode the challenge seed and the transmitted response $z_1$, then
        recompute the public binary challenge $c$.
  \item For every $j$ with $c[j]=1$, update
        \begin{equation}
          C_{p,d}\mathrel{+}=z_{1,0}[j]\,
                      z_{1,p}[j+d]_{\rm neg}
          \quad\text{for all }p,d,
        \label{eq:update}
        \end{equation}
        and increment $M$.
\end{enumerate}

At a checkpoint, set
\begin{equation}
             \widehat s_{1,p}[d]
               =\round{\frac{C_{p,d}}{2M}},
\label{eq:estimator}
\end{equation}
clamped to the public encoding interval for a secret coefficient.  The scale
$2$ is fixed by Equation~\eqref{eq:rankone}; it is not fitted from the unknown
key.  The estimator stores $O(\ell n)$ words.  The straightforward update
cost is $O(Q\tau\ell n)$ for $Q$ signatures.  FFT-based batched correlation
could reduce this cost, but was unnecessary for the reported laptop runs.

\subsection{Public completion and exact validation}

Once $\widehat s_1$ is available, expand $a_{\rm gen}$ and $A_{\rm gen}$ from
the public seed and compute
\begin{equation}
 \widehat s'_2=\operatorname{center}_q
       \left(b-a_{\rm gen}-A_{\rm gen}\widehat s_1\right).
\label{eq:errorcompletion}
\end{equation}
This is public computation.  A candidate is accepted only if

\begin{itemize}[leftmargin=1.5em]
  \item every coefficient of $\widehat s_1$ and $\widehat s'_2$ is inside
        the submitted secret-key encoding range;
  \item $\Stretch(1,\widehat s_1,\widehat s'_2)$ lies inside the submitted
        two-sided key-generation norm window; and
  \item Equation~\eqref{eq:pkrelation} holds exactly.
\end{itemize}

These checks use no original secret information.  In the reported
experiments, public validation first succeeded at the same checkpoint at which
all coefficients of $s_1$ matched the experimental ground truth.  An accepted
fresh-message signature produced by the completed key is the final end-to-end
validation.

\subsection{Equivalent-key forgery}

The serialized Shuttle secret key contains the public seed and $b$, an
auxiliary signing seed $K$, the public-key hash, and $(s_1,s'_2)$.  After
Equations~\eqref{eq:estimator} and \eqref{eq:errorcompletion}, choose any new
$K$, recompute the public-key hash, and encode these values in the documented
secret-key format.  This key need not equal the original byte string.  It
satisfies the same public relation and all signing bounds, so the public
signing algorithm can sign an arbitrary message.  Verification uses only the
original public key.

The demonstrated attack operates in the EUF-CMA setting: it obtains valid
signatures on distinct messages and subsequently produces an accepted
signature on a fresh message that was not queried before.

\section{Parameters and experiments}
\label{sec:experiments}

\subsection{Submitted parameter sets}

Table~\ref{tab:params} lists the quantities relevant to the attack.  All three
sets use $\ell=3$ and Gaussian width $r=825$.

\begin{table}[ht]
\centering
\caption{Relevant parameters of the submitted Shuttle instances.}
\label{tab:params}
\small
\begin{tabular}{@{}lrrrrrrrr@{}}
\toprule
Instance & $n$ & $q$ & $m$ & $\tau$ & $\alpha_1$ & $\alpha_s$ & $B_s$ & Sig. bytes\\
\midrule
Shuttle-128 & 256  & 15361 & 3 & 42  & 90  & 10 & 9  & 1183\\
Shuttle-256 & 512  & 61441 & 2 & 58  & 135 & 5  & 10 & 2417\\
Shuttle-512 & 1024 & 59393 & 2 & 115 & 144 & 3  & 10 & 5001\\
\bottomrule
\end{tabular}
\end{table}

\subsection{Methodology}

We used the portable C reference implementation from the public NGCC archive,
whose SHA-256 digest is
\begin{center}
\scriptsize\texttt{09253fd33c2215098177991db3375cb94e4e4a73a65b019783edabefa5bd8a62}.
\end{center}
The scheme sources were not modified.  They were compiled with GCC 13.3.0 at
\texttt{-O3} and \texttt{-march=native}.  The attack ran as one process and did not use
worker threads.  The machine had an Intel Core i5-6500 at 3.2 GHz, four
physical cores, and 16 GiB of RAM; CPU times therefore describe one ordinary
core and are more reproducible than wall times under concurrent system load.

For each query we signed a distinct 32-byte message using the implementation's
explicit-randomness API, verified the result, decoded it through the public
signature decoder, and applied Equation~\eqref{eq:update}.  Checkpoints were
taken every 25,000 signatures.  Recovery used the fixed divisor $2M$.  A
least-squares scale computed using the experimental ground truth was printed
only as a diagnostic after the estimate; it was never an attack input.

Success required all public completion checks and an accepted signature on
the fresh message
\begin{center}
\texttt{fresh message signed with the statistically recovered key}.
\end{center}
Every honest query and every reported forgery was checked by the unmodified
submitted verifier.

\subsection{Results}

Table~\ref{tab:results} gives the first checkpoint at which the direct rounded
estimate passed public validation.  A diagnostic regression slope between
average covariance and the known experimental secret stayed between 1.95 and
2.03 during the runs.  Its proximity to $2$ supports the scale predicted by
Equation~\eqref{eq:observable}; the actual attack always used the fixed value
$2$ and never the diagnostic fit.

\begin{table}[ht]
\centering
\caption{End-to-end equivalent-key recovery and fresh-message forgery.}
\label{tab:results}
\small
\begin{tabular}{@{}llrrrl@{}}
\toprule
Instance & Key tag & Signatures & Observations & CPU s & Forgery\\
\midrule
Shuttle-128 & \texttt{42} & 200,000 & 8,400,000 & 413.175 & accept\\
Shuttle-128 & \texttt{43} & 150,000 & 6,300,000 & 310.879 & accept\\
Shuttle-128 & \texttt{44} & 225,000 & 9,450,000 & 467.112 & accept\\
Shuttle-128 & \texttt{45} & 175,000 & 7,350,000 & 362.885 & accept\\
Shuttle-256 & \texttt{42} & 225,000 & 13,050,000 & 739.804 & accept\\
Shuttle-512 & \texttt{42} & 325,000 & 37,375,000 & 2208.998 & accept\\
\bottomrule
\end{tabular}
\end{table}

The four Shuttle-128 trials used independent deterministic key seeds.  Direct
rounding recovered all $3n=768$ coefficients and public completion recovered
all 768 or 512 error coefficients, as appropriate.  The Shuttle-256 and
Shuttle-512 trials recovered all 1536 and 3072 coefficients of $s_1$ and all
1024 and 2048 coefficients of $s'_2$, respectively.  The generated forgery
lengths were exactly 1183, 2417, and 5001 bytes.

Convergence was smooth rather than a rare event.  For the Shuttle-512 trial,
the exact secret-coefficient counts at 25,000-signature intervals were
2189, 2652, 2856, 2965, 3016, 3045, 3061, 3066, 3068, 3069, 3070, 3070, and
3072 out of 3072.  Incomplete checkpoints failed the public validation
test.  The fitted scale stayed within a few percent of $2$ throughout.

The attack state itself consists principally of $3n$ 64-bit accumulators:
6 KiB, 12 KiB, and 24 KiB for the three levels.  The measured peak resident
set, including the submitted signer, verifier, decoding buffers, and process
runtime, stayed below 2.5 MB.  Transcript storage is unnecessary.

\subsection{Corrected-sign control}
\label{sec:corrected-control}

To isolate the cause, we made one experimental change to a separate copy of
the Shuttle-128 reference implementation: the final transition update was
changed from $z\leftarrow z+Bv$ to $z\leftarrow z-Bv$.  No attack or verifier
code was changed.  All 25,000 generated signatures verified, but the
ground-truth diagnostic slope fell from approximately $2$ to $0.03799$ and
Equation~\eqref{eq:estimator} recovered only 348 of 768 coefficients.  The
candidate failed public validation.  This control is consistent with the
corrected transition suppressing the rank-one covariance term and directly
links the observed key recovery to the sign reversal.

\section{Scope, cause, and repair}
\label{sec:scope}

The reversed branch direction is present in both the NGCC specification
pseudocode and the submitted reference implementation.  The attack is not
caused by the placeholder SM3-based XOF, floating-point rounding, compiler
behavior, timing leakage, or malformed decoding.  Replacing the XOF with a
random oracle leaves the covariance mechanism intact.

The minimal conceptual repair is to apply the interval event to the opposite
shift: after the sign normalization used by the submitted algorithm, update
the state by $y-Bv$ instead of $y+Bv$.  This matches the definition of
$p_v(y)$ and the transition in \cite{MingEtAl}.  Changing that sign changes
known-answer vectors and requires a fresh correctness, distribution, and
constant-time review; implementers should not treat it as an isolated source
patch without regenerating and independently validating the submission.

The attack does not claim a break of the correctly oriented construction in
the separately published ePrint paper.  That construction has a different
public parameterization and states the branch direction consistently.  Our
result is specific to the public NGCC specification and the code submitted
with it.

\section{Disclosure and prior-art check}

Before this manuscript was prepared, we searched the public NGCC report index,
the pqc-forum web archive through general web search, the IACR ePrint index,
and a contemporary roundup of NGCC findings \cite{NGCCReports,Roundup}.  On 25
September 2026 the NGCC index listed Shuttle (sign-23) as ``No report,'' and
the roundup did not mention Shuttle.  We found the authors' ePrint
\cite{MingEtAl}, which is directly relevant because it specifies the correct
branch direction, but it does not report the reversed NGCC implementation or
this key-recovery attack.  This is a time-bounded search, not proof of novelty;
human review remains pending.

\section{Conclusion}

In the NGCC submission, the interval event associated in the specification
with the $y-v$ branch is applied by Algorithm~22 and the reference
implementation to the $y+v$ update.  In the idealized one-transition analysis,
this sign reversal replaces covariance cancellation with an additional
$2vv^T$ secret-dependent term.  The known constant component of Shuttle's
secret makes the term directly observable in ordinary signatures.  A simple
streaming correlation recovers the complete signing secret, the public key
determines the omitted error, and the resulting equivalent key forges fresh
messages for every submitted parameter set.  The largest demonstrated run
uses well below one million signatures, one CPU core, and negligible attack
memory.

\begin{thebibliography}{9}

\bibitem{Shuttle}
L. Chen, Y. Ming, J. Zhang, C. Zhang, Z. Liu, G. Tang, S. Gao, P. Chen,
Y. Sun, T. Wang, M. Mei, and J. Hou,
\emph{Shuttle: A Lattice Based Digital Signature Scheme---Algorithm
Specification and Supporting Documentation}, NGCC submission, 30 June 2026.

\bibitem{MingEtAl}
Y. Ming, J. Zhang, Z. Liu, G. Tang, P. Chen, Y. Sun, S. Gao, C. Zhang, and
L. Chen,
\emph{Towards Practical Iterative Rejection Sampling: A Compact and Efficient
Signature over Module Lattices}, IACR Cryptology ePrint Archive, Report
2026/1991, 2026.  \url{https://eprint.iacr.org/2026/1991}

\bibitem{Gaertner}
J. G\"artner,
\emph{Compact Lattice Signatures via Iterative Rejection Sampling},
Advances in Cryptology--CRYPTO 2025, LNCS 16000, pp. 514--547, 2025.

\bibitem{NGCCReports}
M.-J. O. Saarinen,
\emph{NGCC Reports}, public report index, accessed 25 September 2026.
\url{https://ngcc.dev/reports/index.html}

\bibitem{Roundup}
M. Ivezic,
\emph{China's Next-Generation Crypto Candidates Drew 104 Public Findings in
Three Days}, PostQuantum.com, 23 September 2026.
\url{https://postquantum.com/security-pqc/china-ngcc-candidate-flaws/}

\end{thebibliography}

\end{document}
